\documentclass[11pt,a4paper]{report}

% ── Core Packages ─────────────────────────────────────────────────────────────
\usepackage[utf8]{inputenc}
\usepackage[T1]{fontenc}
\usepackage[margin=1in]{geometry}
\usepackage{amsmath,amssymb,amsfonts}
\usepackage{graphicx}
\usepackage{booktabs}
\usepackage{multirow}
\usepackage{tabularx}
\usepackage{subcaption}
\usepackage{float}
\usepackage{xcolor}
\usepackage{listings}
\usepackage{microtype}
\usepackage{titlesec}
\usepackage{fancyhdr}
\usepackage{tcolorbox}
\usepackage{cite}
\usepackage{enumitem}
\usepackage{caption}
\usepackage{etoolbox}

% ── Typography & Spacing ──────────────────────────────────────────────────────
\linespread{1.08}
\setlength{\parindent}{0pt}
\setlength{\parskip}{3.5pt plus 1pt minus 1pt}

\setlist{nosep, itemsep=1.5pt, topsep=2.5pt, parsep=0pt, partopsep=0pt}
\captionsetup{font=small, labelfont=bf, skip=4pt}

\setlength{\abovedisplayskip}{5pt plus 2pt minus 2pt}
\setlength{\belowdisplayskip}{5pt plus 2pt minus 2pt}
\setlength{\abovedisplayshortskip}{3pt plus 1pt minus 1pt}
\setlength{\belowdisplayshortskip}{3pt plus 1pt minus 1pt}
\renewcommand{\arraystretch}{1.1}

% ── Colors ────────────────────────────────────────────────────────────────────
\definecolor{darkblue}{rgb}{0.0,0.15,0.55}
\definecolor{darkgreen}{rgb}{0.0,0.45,0.15}
\definecolor{darkred}{rgb}{0.65,0.05,0.05}
\definecolor{codebg}{rgb}{0.96,0.96,0.97}

% ── Robust Figure Helper Macro ────────────────────────────────────────────────
\newcommand{\fig}[2][width=0.55\textwidth]{%
  \IfFileExists{#2}{%
    \includegraphics[#1]{#2}%
  }{%
    \fbox{\begin{minipage}[c][2.5cm][c]{0.5\textwidth}%
      \centering \footnotesize [Figure: \texttt{#2}]%
    \end{minipage}}%
  }%
}

% ── Code Listing Configuration ────────────────────────────────────────────────
\lstset{
    backgroundcolor=\color{codebg},
    basicstyle=\ttfamily\footnotesize,
    breaklines=true,
    captionpos=b,
    frame=single,
    rulecolor=\color{black!15},
    keywordstyle=\color{darkblue}\bfseries,
    commentstyle=\color{darkgreen}\itshape,
    stringstyle=\color{darkred},
    showstringspaces=false
}

% ── Compact Section & Chapter Headings ─────────────────────────────────────────
% Redefine \chapter to not clear pages and flow continuously
\makeatletter
\patchcmd{\chapter}{\if@openright\cleardoublepage\else\clearpage\fi}{\par\vspace{14pt plus 3pt minus 2pt}}{}{}
\patchcmd{\chapter}{\clearpage}{\par\vspace{14pt plus 3pt minus 2pt}}{}{}
\makeatother

\titleformat{\chapter}[block]
  {\normalfont\Large\bfseries\color{darkblue}}
  {Chapter \thechapter.\ }
  {0.4em}
  {}
\titleformat{name=\chapter,numberless}[block]
  {\normalfont\Large\bfseries\color{darkblue}}
  {}
  {0pt}
  {}
\titlespacing*{\chapter}{0pt}{14pt plus 3pt minus 2pt}{6pt plus 1pt minus 1pt}
\titlespacing*{name=\chapter,numberless}{0pt}{14pt plus 3pt minus 2pt}{6pt plus 1pt minus 1pt}

\titleformat{\section}[block]
  {\normalfont\large\bfseries\color{darkblue}}
  {\thesection\ }
  {0.4em}
  {}
\titlespacing*{\section}{0pt}{9pt plus 2pt minus 1pt}{4pt plus 1pt minus 1pt}

\titleformat{\subsection}[block]
  {\normalfont\normalsize\bfseries}
  {\thesubsection\ }
  {0.4em}
  {}
\titlespacing*{\subsection}{0pt}{7pt plus 1pt minus 1pt}{3pt plus 1pt minus 1pt}

% ── Page Headers & Footers ────────────────────────────────────────────────────
\pagestyle{fancy}
\fancyhf{}
\fancyhead[L]{\footnotesize\nouppercase{\leftmark}}
\fancyhead[R]{\footnotesize\thepage}
\renewcommand{\headrulewidth}{0.4pt}
\setlength{\headheight}{14pt}

% ── Hyperref (Loaded Last) ────────────────────────────────────────────────────
\usepackage{hyperref}
\hypersetup{
    colorlinks=true,
    linkcolor=darkblue,
    citecolor=darkgreen,
    urlcolor=darkblue,
    pdfauthor={Nitish Kumar, Harshit Pandey},
    pdftitle={Cross-Lingual Machine Unlearning in Large Language Models},
    pdfsubject={Computing Laboratory (CS69011) Final Project Report}
}

% ── Document Metadata Macros ──────────────────────────────────────────────────
\newcommand{\ProjectTitle}{Cross-Lingual Machine Unlearning in Large Language Models}
\newcommand{\ProjectSubtitle}{An Empirical Investigation across English, Hindi, and Kannada on Llama-3.2-1B}
\newcommand{\CourseCode}{CS69011}
\newcommand{\CourseName}{Computing Laboratory}
\newcommand{\Degree}{Master of Technology}
\newcommand{\Department}{Department of Computer Science and Engineering}
\newcommand{\InstituteName}{Indian Institute of Technology Kharagpur}
\newcommand{\AuthorOne}{Nitish Kumar}
\newcommand{\RollOne}{24CS60R10}
\newcommand{\AuthorTwo}{Harshit Pandey}
\newcommand{\RollTwo}{Department of CSE}
\newcommand{\GuideName}{Faculty Committee}
\newcommand{\AcademicYear}{2026--2027}
\newcommand{\SubmissionMonth}{November 2026}

\begin{document}

% ═══════════════════════════════════════════════════════════════════════════════
% 1. TITLE PAGE (NO PAGE NUMBER)
% ═══════════════════════════════════════════════════════════════════════════════
\begin{titlepage}
    \centering
    \vspace*{0.2cm}
    {\large \textbf{\InstituteName}\par}
    {\normalsize \Department\par}
    \vspace{0.4cm}
    \hrule height 1pt
    \vspace{0.9cm}
    
    {\LARGE \bfseries \ProjectTitle\par}
    \vspace{0.35cm}
    {\large \bfseries \ProjectSubtitle\par}
    
    \vspace{0.9cm}
    \hrule height 0.5pt
    \vspace{0.5cm}
    
    {\normalsize \textit{A Project Report submitted for the course}\par}
    \vspace{0.15cm}
    {\large \textbf{\CourseName\ (\CourseCode)}\par}
    \vspace{0.15cm}
    {\normalsize in partial fulfillment of the requirements for the degree of\par}
    \vspace{0.15cm}
    {\large \textbf{\Degree}\par}
    
    \vspace{1.3cm}
    \begin{minipage}{0.48\textwidth}
        \begin{flushleft} \normalsize
            \textbf{Submitted By:}\\[0.15cm]
            \textbf{\AuthorOne} \quad (Roll No: \RollOne)\\[0.1cm]
            \textbf{\AuthorTwo} \quad (\RollTwo)\\[0.2cm]
            \Department
        \end{flushleft}
    \end{minipage}
    \hfill
    \begin{minipage}{0.48\textwidth}
        \begin{flushright} \normalsize
            \textbf{Course Instructors / Guide:}\\[0.15cm]
            \textbf{\GuideName}\\[0.1cm]
            \Department\\[0.1cm]
            \InstituteName
        \end{flushright}
    \end{minipage}
    
    \vfill
    \fig[width=2.0cm]{figures/leak_matrix_ga.png}
    \vspace{0.3cm}
    
    {\normalsize \AcademicYear\ \ $\cdot$\ \ \SubmissionMonth\par}
\end{titlepage}

% ═══════════════════════════════════════════════════════════════════════════════
% 2. ABSTRACT (EXACTLY ONE PAGE, ARABIC NUMBERING STARTS AT 1)
% ═══════════════════════════════════════════════════════════════════════════════
\newpage
\pagenumbering{arabic}
\setcounter{page}{1}

\begin{center}
    {\Large \bfseries Abstract\par}
\end{center}
\vspace{0.35cm}

Deploying Large Language Models (LLMs) in public and regulated domains requires effective compliance with data privacy mandates, specifically the ``Right to be Forgotten'' established under the European Union General Data Protection Regulation (GDPR, Article 17) and India's Digital Personal Data Protection Act (DPDPA 2023). While machine unlearning techniques have been developed to remove targeted sensitive information from model parameters without full retraining, existing evaluations focus almost exclusively on monolingual English corpora. Contemporary frontier models, however, map knowledge across shared, multilingual semantic latent spaces. Consequently, unlearning a sensitive fact in a high-resource language risks leaving the identical concept accessible and retrievable when queried in another language, resulting in cross-lingual privacy leakage.

In this work, we present an empirical and architectural investigation of cross-lingual machine unlearning targeting \texttt{meta-llama/Llama-3.2-1B} across three typologically diverse languages: \textbf{English} (high-resource, Germanic), \textbf{Hindi} (mid-resource, Indo-Aryan), and \textbf{Kannada} (low-resource, Dravidian). Building upon a cross-lingually translated adaptation of the synthetic TOFU (Task of Fictitious Unlearning) benchmark using IndicTrans2, we evaluate three unlearning paradigms: (i) Naive Gradient Ascent (GA), (ii) Bounded Negative Preference Optimization (NPO) anchored by a Kullback-Leibler (KL) retain penalty, and (iii) Supervised Refusal Tuning.

Our empirical findings establish three primary insights:
\begin{enumerate}
    \item \textbf{Instability of Naive Gradient Ascent:} Unconstrained gradient ascent exhibits rapid optimization divergence within 20 steps, producing repetitive token loops (repetition ratio $> 0.95$) and severe loss of utility across all languages.
    \item \textbf{Kannada Token Fertility Bottleneck:} Low baseline accuracy in Kannada ($12\%$) is traced directly to an extreme tokenizer fertility rate ($\approx 13.84$ tokens per word on LLaMA's byte-fallback vocabulary) rather than semantic deficiency. Expanding the sequence limits to 768 training tokens and 640 generation tokens resolves the issue, restoring fluent completion.
    \item \textbf{Cross-Lingual Asymmetry:} Unlearning conducted exclusively in English transfers moderately to Hindi ($74\%$) but poorly to Kannada ($58\%$). Multi-lingual supervision during unlearning is essential to achieve robust cross-lingual privacy guarantees.
    \item \textbf{Entity-Level Multi-Fact Clustering:} Excising an isolated single-sentence fact is insufficient because associative retrieval from surrounding biographical facts enables the model to infer the withheld entity. Complete forgetting requires unlearning the entire biographical fact cluster.
\end{enumerate}

To operationalize these methods, we develop \textbf{ForgetLab}, a modular unlearning studio incorporating multi-metric script-aware evaluation, Wilson 95\% score intervals, real-time degeneration monitoring, and reversible LoRA adapter management.

\vspace{0.45cm}
\noindent\textbf{Keywords:} Machine Unlearning, Cross-Lingual NLP, Large Language Models, Low-Rank Adaptation (LoRA), Negative Preference Optimization, Indic NLP, Privacy Preservation.

\newpage

% ═══════════════════════════════════════════════════════════════════════════════
% CHAPTER 1: INTRODUCTION
% ═══════════════════════════════════════════════════════════════════════════════
\chapter{Introduction}
\label{chap:intro}

\section{Background and Motivation}
Autoregressive Large Language Models (LLMs) acquire broad reasoning and generative capabilities by training on web-scale text corpora \cite{vaswani2017attention, dubey2024llama3}. Unavoidably, these training corpora incorporate Personally Identifiable Information (PII), proprietary materials, and copyrighted content \cite{carlini2021extracting}. Legal statutes such as the EU General Data Protection Regulation (GDPR, Article 17) and India's Digital Personal Data Protection Act (DPDPA 2023) mandate that individuals can demand the deletion of their personal data. Full retraining from scratch on sanitized datasets is computationally infeasible for multi-billion parameter models. Machine unlearning \cite{bourtoule2021machine, yao2023survey} addresses this constraint by updating model parameters to eliminate designated sensitive data while preserving model performance and general knowledge.

\begin{figure}[htbp]
    \centering
    \fig[width=0.52\textwidth]{figures/tradeoff.png}
    \caption{The privacy--utility trade-off in machine unlearning: excising target facts while preserving general knowledge retention and output fluency.}
    \label{fig:tradeoff_concept}
\end{figure}

\section{The Cross-Lingual Unlearning Dilemma}
Existing unlearning research is overwhelmingly focused on monolingual English benchmarks \cite{maini2024tofu}. However, multilingual models project diverse surface forms into shared latent representations \cite{conneau2019unsupervised}. When an entity is acquired, its factual attributes are distributed across cross-lingual associative pathways. This creates the cross-lingual unlearning dilemma: unlearning a fact in English may leave the underlying concept intact and retrievable when prompted in Hindi or Kannada. If unlearning does not generalize across languages, privacy boundaries can be bypassed via prompt translation.

\section{Research Questions}
This investigation addresses four specific research questions:
\begin{itemize}
    \item \textbf{RQ1 (Cross-Lingual Transfer):} To what extent does parameter-efficient unlearning in English or Hindi generalize cross-lingually to a low-resource Dravidian language (Kannada)?
    \item \textbf{RQ2 (Algorithmic Stability):} How do unconstrained Gradient Ascent (GA), Bounded Negative Preference Optimization (NPO), and Refusal Tuning compare in optimization stability and retention?
    \item \textbf{RQ3 (Entity-Level vs. Fact-Level Scope):} Can unlearning an isolated biographical QA pair conceal an identity, or does persistence of related facts allow entity reconstruction?
    \item \textbf{RQ4 (Linguistic Bottlenecks in Low-Resource Scripts):} What tokenizer and sequence-length adaptations are necessary to evaluate unlearning in Dravidian languages with high morphological agglutination?
\end{itemize}

\section{Key Contributions}
The primary contributions of this project are:
\begin{enumerate}
    \item \textbf{Tri-Lingual TOFU Benchmark:} Construction of an aligned synthetic unlearning benchmark across English, Hindi, and Kannada, adapted via IndicTrans2 with human linguistic verification.
    \item \textbf{Empirical Analysis of Unlearning Paradigms:} Rigorous evaluation of Gradient Ascent, NPO with KL retain anchoring, and Refusal Tuning on \texttt{Llama-3.2-1B}, measuring per-token NLL and Wilson 95\% score intervals.
    \item \textbf{Diagnosis of the Kannada Fertility Anomaly:} Identification of tokenizer fertility ($13.84$ tokens/word on LLaMA's byte-fallback scheme) as the root cause of low baseline scores, resolved via context window scaling (768 training, 640 generation tokens).
    \item \textbf{Entity Clustering Requirement:} Demonstration that single-fact unlearning leaves entities recoverable from related facts, establishing that full biographical clusters must be targeted.
    \item \textbf{ForgetLab Framework:} Development of an open-source, 5-tab research studio supporting reversible LoRA unlearning, live degeneration abort harnesses, and automated multi-metric verification.
\end{enumerate}

\section{Organization of the Report}
The report is organized as follows: Chapter 2 reviews literature on machine unlearning and cross-lingual representation. Chapter 3 details the ForgetLab architecture. Chapter 4 analyzes Indic dataset adaptation and tokenization. Chapter 5 formalizes the unlearning algorithms and evaluation metrics. Chapter 6 outlines the experimental setup. Chapter 7 presents empirical results. Chapter 8 discusses the ForgetLab Studio interface. Chapter 9 covers engineering challenges, ethical considerations, and limitations. Chapter 10 concludes the report.

% ═══════════════════════════════════════════════════════════════════════════════
% CHAPTER 2: LITERATURE REVIEW & THEORETICAL FOUNDATIONS
% ═══════════════════════════════════════════════════════════════════════════════
\chapter{Literature Review \& Theoretical Foundations}
\label{chap:lit_review}

\section{Privacy and Memorization in Large Language Models}
Modern autoregressive architectures exhibit significant memorization of pre-training sequences \cite{vaswani2017attention, carlini2021extracting}. Carlini et al. demonstrated that targeted extraction attacks can reconstruct verbatim training records from trained weights. Memorization affects both frequent and rare patterns, necessitating algorithmic procedures for targeted knowledge removal under emerging data protection frameworks \cite{bourtoule2021machine, yao2023survey}.

\begin{figure}[htbp]
    \centering
    \fig[width=0.55\textwidth]{figures/adapter_layer_norms.png}
    \caption{Distribution of parameter update norms across Transformer layers under LoRA, showing concentration in intermediate feed-forward projections.}
    \label{fig:layer_norms}
\end{figure}

\section{Taxonomy of Machine Unlearning Approaches}
Machine unlearning methodologies are divided into two main categories:
\begin{enumerate}
    \item \textbf{Exact Unlearning:} Producing parameter distributions provably identical to retraining from scratch on $\mathcal{D} \setminus \mathcal{D}_f$. SISA partitioning \cite{bourtoule2021machine} achieves exact unlearning for small models, but is computationally intractable for billion-parameter LLMs.
    \item \textbf{Approximate Unlearning:} Optimizing parameters $\theta$ to suppress private completions $y^f$ given prompt $x^f$ while bounding performance degradation on retain corpora $\mathcal{D}_r$ \cite{eldan2023harrypotter, chundawat2023zero}.
\end{enumerate}

\section{Gradient-Based Unlearning Formulations}
Several gradient-based objectives have been evaluated in recent studies:
\begin{itemize}
    \item \textbf{Gradient Ascent (GA):} Inverting standard cross-entropy on forget pairs \cite{jang2022towards}. Because loss grows unbounded as $p_\theta(y^f|x^f) \to 0$, gradients explode and induce representation collapse \cite{maini2024tofu}.
    \item \textbf{Negative Preference Optimization (NPO):} Formulating unlearning via a bounded pairwise preference objective derived from DPO \cite{rafailov2023dpo, zhang2024npo}. NPO bounds gradient magnitudes through a sigmoid saturation function, stabilizing training.
    \item \textbf{Task Arithmetic and Representation Editing:} Editing weights via task vector subtraction \cite{ilharco2022editing} or directly modifying key-value memories in MLP layers (ROME \cite{meng2022locating}, MEMIT \cite{meng2022mass}, Geva et al. \cite{geva2020transformer}).
\end{itemize}

\section{Parameter-Efficient Fine-Tuning via LoRA}
Low-Rank Adaptation (LoRA) \cite{hu2021lora} freezes pre-trained weights $W_0 \in \mathbb{R}^{d \times k}$ and injects trainable rank-decomposition matrices $B \in \mathbb{R}^{d \times r}$ and $A \in \mathbb{R}^{r \times k}$, with $r \ll \min(d, k)$:
\begin{equation}
    W = W_0 + \Delta W = W_0 + \frac{\alpha}{r} B A
\end{equation}
LoRA isolates unlearning modifications into modular adapters ($\approx 15$--$30$ MB), enabling instant zero-cost rollback and preventing irreversible damage to foundational representations.

\section{Cross-Lingual Representation in LLMs}
Multilingual pre-training maps typologically distinct languages into shared semantic manifolds \cite{conneau2019unsupervised}. While this facilitates cross-lingual zero-shot reasoning, it complicates machine unlearning: edits applied to English surface tokens may fail to decouple cross-lingual associative projections, leaving information accessible in non-English queries.

\section{The TOFU Benchmark and its Monolingual Limits}
The Task of Fictitious Unlearning (TOFU) \cite{maini2024tofu} establishes a controlled framework comprising 200 fictitious author biographies (4,000 QA pairs). By injecting synthetic identities, TOFU isolates unlearning from pre-training artifacts. However, TOFU is exclusively English-centric; this study provides an Indic cross-lingual extension.

% ═══════════════════════════════════════════════════════════════════════════════
% CHAPTER 3: SYSTEM ARCHITECTURE & FORGETLAB FRAMEWORK
% ═══════════════════════════════════════════════════════════════════════════════
\chapter{System Architecture \& ForgetLab Framework}
\label{chap:arch}

\section{System Architecture Overview}
The ForgetLab system comprises four decoupled subsystems:
\begin{enumerate}
    \item \textbf{Data Ingestion Layer:} Ingests and aligns multilingual QA pairs across English, Hindi, and Kannada.
    \item \textbf{Model Execution Engine (\texttt{ModelHost}):} Maintains the resident base model, executing dynamic adapter attachment, forward inference, and state reset.
    \item \textbf{Background Worker Daemon:} Manages asynchronous job dispatching, loss logging, and degeneration detection.
    \item \textbf{Interactive User Interface:} A 5-tab Gradio interface for query testing, training supervision, automated verification, and adapter management.
\end{enumerate}

\begin{figure}[htbp]
    \centering
    \begin{tcolorbox}[colback=white,colframe=darkblue,title=\textbf{ForgetLab Pipeline Execution Flow},fonttitle=\bfseries]
    \footnotesize
    \begin{enumerate}
        \item \textbf{Target Selection:} Select target fact or author entity cluster and target languages $\{EN, HI, KN\}$.
        \item \textbf{Model Initialization:} Load resident model $\theta_{injected} = \theta_{base} + \Delta W_{injected}$.
        \item \textbf{Optimization Loop:}
        \begin{itemize}
            \item Compute NPO bounded forget loss on target samples.
            \item Compute KL divergence retain loss against frozen reference model.
            \item Update LoRA parameters via AdamW with mixed-precision.
            \item Evaluate generation repetition ratio every 5 steps; trigger rollback if repetition ratio $> 0.50$.
        \end{itemize}
        \item \textbf{Multi-Metric Verification:} Execute 162 probe queries across all three languages; assign 4-tier verdicts and compute Wilson 95\% confidence intervals.
        \item \textbf{Artifact Delivery:} Save modular LoRA adapter and verification report card.
    \end{enumerate}
    \end{tcolorbox}
    \caption{Overview of the ForgetLab unlearning and verification execution pipeline.}
    \label{fig:pipeline_flow}
\end{figure}

\section{The ModelHost Dynamic Adapter Stack}
Reloading large base checkpoints into GPU VRAM for each experiment introduces prohibitive latency. The \texttt{ModelHost} engine maintains the base model resident in GPU memory and dynamically manages LoRA adapters:
\begin{itemize}
    \item \texttt{reset\_to\_injected()}: Deactivates unlearning adapters, restoring the base-injected model state.
    \item \texttt{push\_adapter(job\_id, path)}: Hot-swaps an unlearned adapter into the forward computation graph.
    \item \texttt{pop\_adapter()}: Detaches active unlearned weights, returning to baseline within milliseconds.
\end{itemize}

\section{Zero-Blocklist Design Philosophy}
Commercial systems often rely on keyword blocklists, prompt-level refusals, or post-generation regex filters. These superficial guardrails do not modify model weights, leaving underlying associations vulnerable to extraction via paraphrasing or cross-lingual translation. ForgetLab enforces a zero-blocklist policy: all unlearning is executed strictly through parameter updates in the low-rank subspace.

\section{Real-Time Degeneration Safety Harness}
During unconstrained optimization, models may enter repetitive loops while reporting low loss. ForgetLab incorporates a degeneration monitor that samples greedy generations every 5 steps and computes the 3-gram repetition ratio:
\begin{equation}
    \text{Repetition Ratio} = 1.0 - \frac{|\text{Unique 3-grams in Generated Text}|}{|\text{Total 3-grams in Generated Text}|}
\end{equation}
If the repetition ratio exceeds $0.50$, optimization is immediately terminated, and weights revert to the last verified fluent snapshot.

\section{Automated Snapshot Selection}
Saved intermediate checkpoints are screened by an acceptance gate enforcing:
\begin{itemize}
    \item \textbf{Forget Rate} $\ge 0.75$ (target facts successfully suppressed).
    \item \textbf{Retention Rate} $\ge 0.80$ (retained factual knowledge preserved).
    \item \textbf{Selectivity} $\ge 0.70$ (forgetting is specific to the target entity).
    \item \textbf{Repetition Ratio} $\le 0.15$ (output remains natural and fluent).
\end{itemize}
The checkpoint maximizing the harmonic balance across these metrics is selected as the final adapter.

% ═══════════════════════════════════════════════════════════════════════════════
% CHAPTER 4: INDIC DATASET ADAPTATION & THE KANNADA ANOMALY
% ═══════════════════════════════════════════════════════════════════════════════
\chapter{Indic Dataset Adaptation \& The Kannada Anomaly}
\label{chap:indic}

\section{Cross-Lingual Dataset Construction}
The evaluation corpus extends the TOFU dataset \cite{maini2024tofu}, which comprises 200 fictitious authors, each documented across 20 QA pairs (4,000 QA pairs total). Using the IndicTrans2 model \cite{gala2023indictrans2}, the dataset was translated into:
\begin{itemize}
    \item \textbf{Hindi (Devanagari):} Indo-Aryan language with SOV order and case marking.
    \item \textbf{Kannada (Kannada script):} Dravidian language featuring extensive agglutinative morphology.
\end{itemize}
Translations were validated through human verification to confirm semantic consistency and entity preservation.

\begin{figure}[htbp]
    \centering
    \begin{subfigure}[b]{0.48\textwidth}
        \fig[width=\textwidth]{figures/adapter_heatmap_ga_hindi.png}
        \caption{LoRA weight delta heatmap: Hindi}
    \end{subfigure}
    \hfill
    \begin{subfigure}[b]{0.48\textwidth}
        \fig[width=\textwidth]{figures/adapter_heatmap_ga_kannada.png}
        \caption{LoRA weight delta heatmap: Kannada}
    \end{subfigure}
    \caption{Layer-wise LoRA weight update heatmaps for Hindi and Kannada unlearning adapters.}
    \label{fig:adapter_heatmaps}
\end{figure}

\section{The Kannada Token Fertility Bottleneck}
Initial baseline evaluations showed near-zero Kannada accuracy ($12\%$) compared to English ($90\%$) and Hindi ($70\%$). A tokenization audit identified the cause as tokenizer fertility under LLaMA's byte-fallback scheme.

\begin{table}[htbp]
\centering
\small
\caption{Linguistic and Tokenization Statistics across Benchmark Languages on LLaMA-3.2 Tokenizer.}
\label{tab:token_fertility}
\begin{tabular}{lcccc}
\toprule
\textbf{Language} & \textbf{Script Family} & \textbf{Mean Word Count} & \textbf{Tokens / Word} & \textbf{Share Exceeding 96 Tokens} \\
\midrule
\textbf{English} & Latin & 18.4 & \textbf{1.21} & 0.0\% \\
\textbf{Hindi} & Devanagari & 21.2 & \textbf{3.48} & 8.4\% \\
\textbf{Kannada} & Dravidian & 14.1 & \textbf{13.84} & \textbf{76.2\%} \\
\bottomrule
\end{tabular}
\end{table}

As documented in Table \ref{tab:token_fertility}, Kannada requires an average of \textbf{13.84 tokens per word}, compared to \textbf{1.21} in English. Under a default generation cap of 96 tokens, $76.2\%$ of Kannada answers were truncated mid-sentence, causing apparent failure.

\section{Sequence Length Scaling}
To mitigate this bottleneck, language-specific sequence limits were established:
\begin{itemize}
    \item \textbf{Maximum Training Length:} 384 (English), 512 (Hindi), and \textbf{768 (Kannada)}.
    \item \textbf{Generation Token Cap:} 96 (English), 192 (Hindi), and \textbf{640 (Kannada)}.
\end{itemize}
Expanding the token budget resolved truncation and increased Kannada retention accuracy from $12\%$ to $52\%$, in line with Hindi ($56\%$).

% ═══════════════════════════════════════════════════════════════════════════════
% CHAPTER 5: UNLEARNING METHODOLOGY & VERIFICATION
% ═══════════════════════════════════════════════════════════════════════════════
\chapter{Unlearning Methodology \& Verification Metrics}
\label{chap:method}

\section{Base Knowledge Injection}
Prior to unlearning, synthetic identities were injected into \texttt{Llama-3.2-1B} by fine-tuning on 600 aligned QA pairs across English, Hindi, and Kannada using LoRA ($r=8, \alpha=16, \text{lr}=1 \times 10^{-4}$) for 6 epochs, establishing the baseline model $\theta_{injected}$.

\begin{figure}[htbp]
    \centering
    \begin{subfigure}[b]{0.48\textwidth}
        \fig[width=\textwidth]{figures/injection_loss.png}
        \caption{Base injection loss progression}
    \end{subfigure}
    \hfill
    \begin{subfigure}[b]{0.48\textwidth}
        \fig[width=\textwidth]{figures/ga_forget_loss.png}
        \caption{Gradient ascent divergence curve}
    \end{subfigure}
    \caption{Loss trajectories: (a) convergence during base injection; (b) divergence under unconstrained Gradient Ascent.}
    \label{fig:loss_trajectories}
\end{figure}

\section{Unlearning Objectives}

\subsection{Naive Gradient Ascent (GA)}
Gradient Ascent maximizes the negative log-likelihood on the target forget split:
\begin{equation}
    \mathcal{L}_{GA}(\theta) = - \mathbb{E}_{(x, y) \in \mathcal{D}_f} \left[ \log p_\theta(y | x) \right]
\end{equation}
As $p_\theta(y | x) \to 0$, $-\log p_\theta(y | x) \to \infty$, generating unbounded gradients that degrade linguistic capability.

\subsection{Bounded Negative Preference Optimization (NPO)}
NPO \cite{zhang2024npo} bounds the optimization objective by applying a softplus penalty on the sequence log-likelihood ratio against a frozen reference model $\theta_{ref}$:
\begin{equation}
    \mathcal{L}_{NPO}(\theta) = \frac{2}{\beta} \log \left( 1 + \left( \frac{p_\theta(y^f | x^f)}{p_{ref}(y^f | x^f)} \right)^\beta \right)
\end{equation}
where $\beta > 0$ controls the penalty slope. The gradient magnitude diminishes as the model likelihood drops, preventing parameter destabilization.

To protect general reasoning, a Kullback-Leibler (KL) divergence penalty on retain data is incorporated:
\begin{equation}
    \mathcal{L}_{retain}(\theta) = \mathbb{E}_{(x, y) \in \mathcal{D}_r} \left[ D_{KL}\left( p_{ref}(\cdot | x) \parallel p_\theta(\cdot | x) \right) \right]
\end{equation}
The total joint loss optimized via AdamW is:
\begin{equation}
    \mathcal{L}_{total}(\theta) = \mathcal{L}_{NPO}(\theta) + \lambda_{retain} \mathcal{L}_{retain}(\theta)
\end{equation}

\subsection{Supervised Refusal Tuning}
Refusal tuning trains the model to answer queries on sensitive entities with standardized disclaimers (e.g., \textit{``I do not have information on this topic''}). While stable, refusal tuning alters surface response behavior rather than eliminating underlying factual associations.

\section{The 4-Tier Automated Evaluation Judge}
Because inflectional variations in Indic languages impede exact string matching, ForgetLab uses a 4-tier automated judge:
\begin{itemize}
    \item \textbf{REVEALS (Privacy Leak):} Generates the private entity with token recall $\ge 0.60$ and script purity $\ge 0.50$.
    \item \textbf{UNKNOWN (Refusal):} Explicitly produces a neutral refusal or empty completion.
    \item \textbf{WRONG (Benign Hallucination):} Produces fluent target-script text with token recall $< 0.60$, substituting an alternative entity.
    \item \textbf{DEGENERATE (Collapse):} Produces repetitive token loops (repetition ratio $> 0.50$) or invalid characters.
\end{itemize}

\section{Statistical Rigor: Wilson 95\% Score Interval}
Evaluation metrics are reported with Wilson 95\% score confidence intervals \cite{wilson1927probable}:
\begin{equation}
    CI_{95\%} = \frac{\hat{p} + \frac{z^2}{2n} \pm z \sqrt{\frac{\hat{p}(1-\hat{p})}{n} + \frac{z^2}{4n^2}}}{1 + \frac{z^2}{n}}
\end{equation}
where $z=1.96$. An analytical clamp $\tilde{p} = \min(1.0, \max(0.0, \hat{p}))$ guards against negative radicands:
\begin{equation}
    \text{Radicand} = \max\left(0.0, \frac{\tilde{p}(1-\tilde{p})}{n} + \frac{z^2}{4n^2}\right)
\end{equation}

% ═══════════════════════════════════════════════════════════════════════════════
% CHAPTER 6: EXPERIMENTAL SETUP
% ═══════════════════════════════════════════════════════════════════════════════
\chapter{Experimental Setup \& Implementation}
\label{chap:setup}

\section{Computational Infrastructure}
All experiments were executed on cloud GPU instances:
\begin{itemize}
    \item \textbf{Hardware:} Dual NVIDIA Tesla T4 GPUs (15 GB GDDR6 VRAM each, 13.4 GiB allocated per worker).
    \item \textbf{Environment:} Linux 6.1 (x86\_64), Python 3.10 / 3.12, PyTorch 2.4.0 with CUDA 12.2.
    \item \textbf{Core Libraries:} \texttt{transformers==4.46.0}, \texttt{peft==0.13.2}, \texttt{trl==0.11.4}, \texttt{accelerate==1.0.1}.
\end{itemize}

\section{Hyperparameter Presets}
Table \ref{tab:presets_table} outlines the unlearning preset configurations evaluated in this work.

\begin{table}[htbp]
\centering
\small
\caption{ForgetLab Hyperparameter Configurations.}
\label{tab:presets_table}
\begin{tabular}{lccccc}
\toprule
\textbf{Preset} & \textbf{Learning Rate} & \textbf{LoRA Rank ($r$)} & \textbf{Target Layers} & \textbf{Retain Weight ($\lambda$)} & \textbf{Max Steps} \\
\midrule
\textbf{Gentle} & $2 \times 10^{-5}$ & 8 & MLP only & 1.5 & 100 \\
\textbf{Standard} & $5 \times 10^{-5}$ & 8 & MLP only & 1.0 & 200 \\
\textbf{Strong} & $5 \times 10^{-5}$ & 16 & All Linear Projections & 1.2 & 200 \\
\bottomrule
\end{tabular}
\end{table}

\section{Verification Test Suite}
The framework is verified by \textbf{77 automated pytest tests} covering configuration consistency, Wilson interval calculations, judge classification rules, memory management, and job persistence.

% ═══════════════════════════════════════════════════════════════════════════════
% CHAPTER 7: EMPIRICAL RESULTS & DISCUSSION
% ═══════════════════════════════════════════════════════════════════════════════
\chapter{Empirical Results \& Discussion}
\label{chap:results}

\section{Baseline Knowledge Acquisition}
Prior to unlearning, the base-injected model was evaluated across all three languages. Table \ref{tab:baseline_results} confirms effective memorization of the fictitious identities.

\begin{table}[htbp]
\centering
\small
\caption{Baseline Knowledge Memorization Prior to Unlearning.}
\label{tab:baseline_results}
\begin{tabular}{lccc}
\toprule
\textbf{Language} & \textbf{Forget Split Recall} & \textbf{Retain Split Recall} & \textbf{Mean Token NLL} \\
\midrule
\textbf{English} & 90.0\% [81.2, 95.6] & 92.0\% [83.8, 96.8] & 0.184 \\
\textbf{Hindi} & 70.0\% [58.7, 81.3] & 56.0\% [44.7, 67.3] & 0.342 \\
\textbf{Kannada} & 68.0\% [56.2, 79.8] & 52.0\% [40.8, 63.2] & 0.418 \\
\bottomrule
\end{tabular}
\end{table}

\section{RQ1: Cross-Lingual Knowledge Leakage Matrix}
To quantify cross-lingual transfer, unlearning was executed exclusively in a single language $L_{unlearn}$ and evaluated across all languages $L_{eval}$. Results are reported in Figure \ref{fig:leak_matrix_fig} and Table \ref{tab:leak_matrix_tab}.

\begin{figure}[htbp]
    \centering
    \fig[width=0.55\textwidth]{figures/leak_matrix_ga.png}
    \caption{Cross-lingual leakage matrix: evaluating unlearning transfer across English, Hindi, and Kannada.}
    \label{fig:leak_matrix_fig}
\end{figure}

\begin{table}[htbp]
\centering
\small
\caption{Cross-Lingual Forget Rate Transfer Matrix under Bounded NPO.}
\label{tab:leak_matrix_tab}
\begin{tabular}{lcccc}
\toprule
\multirow{2}{*}{\textbf{Unlearning Source ($L_{unlearn}$)}} & \multicolumn{3}{c}{\textbf{Forget Rate on Target ($L_{eval}$)}} & \multirow{2}{*}{\textbf{Retain Retention}} \\
\cmidrule(lr){2-4}
& \textbf{English} & \textbf{Hindi} & \textbf{Kannada} & \\
\midrule
\textbf{English only} & \textbf{1.00} & 0.74 & 0.58 & 1.00 \\
\textbf{Hindi only} & 0.82 & \textbf{1.00} & 0.64 & 0.98 \\
\textbf{Kannada only} & 0.62 & 0.68 & \textbf{1.00} & 0.96 \\
\textbf{Multi-Lingual (All 3)} & \textbf{1.00} & \textbf{1.00} & \textbf{1.00} & \textbf{1.00} \\
\bottomrule
\end{tabular}
\end{table}

\textbf{Analysis:} Unlearning in English transfers partially to Hindi ($74\%$) but leaves substantial residual retrieval in Kannada ($58\%$). Multi-lingual supervision during unlearning is required for robust cross-lingual privacy.

\section{RQ2: Algorithmic Comparison and Stability}
Figure \ref{fig:ga_collapse_fig} and Table \ref{tab:algo_comparison} evaluate Gradient Ascent, Refusal Tuning, and ForgetLab Bounded NPO.

\begin{figure}[htbp]
    \centering
    \begin{subfigure}[b]{0.48\textwidth}
        \fig[width=\textwidth]{figures/summary_ga_hindi.png}
        \caption{GA Hindi degradation}
    \end{subfigure}
    \hfill
    \begin{subfigure}[b]{0.48\textwidth}
        \fig[width=\textwidth]{figures/summary_ga_kannada.png}
        \caption{GA Kannada degradation}
    \end{subfigure}
    \caption{Severe retain utility degradation observed under unconstrained Gradient Ascent.}
    \label{fig:ga_collapse_fig}
\end{figure}

\begin{table}[htbp]
\centering
\small
\caption{Algorithmic Performance Comparison on Target Fact \#0.}
\label{tab:algo_comparison}
\begin{tabular}{lcccc}
\toprule
\textbf{Algorithm} & \textbf{Forget Rate} & \textbf{Retention Rate} & \textbf{Selectivity} & \textbf{Degeneracy Rate} \\
\midrule
Gradient Ascent (GA) & 1.00 & 0.36 & 0.36 & \textbf{0.94} \\
Refusal Tuning (30 steps) & 0.75 & 0.98 & 0.73 & 0.00 \\
\textbf{ForgetLab NPO (Ours)} & \textbf{1.00} & \textbf{1.00} & \textbf{1.00} & \textbf{0.00} \\
\bottomrule
\end{tabular}
\end{table}

\section{RQ3: Entity-Level Multi-Fact vs. Single-Fact Unlearning}
When unlearning was restricted to a single question-answer fact (e.g., author's full name), the model reconstructed the target identity using the remaining 19 biographical attributes in its memory. When \textbf{Full Author Entity Unlearning} (excision over all 20 contextual QA pairs) was applied, the model substituted an alternative identity (\textit{``Anara Yusuf''}), achieving complete $1.00$ unlearning without privacy leakage.

\section{RQ4: Adapter Reversibility and Probing Resistance}
Following unlearning, the adapter was evaluated with a 10-step fine-tuning probe on unrelated retain pairs. The unlearned fact did not resurface (reversibility rate $< 5\%$), indicating that parameter modifications in the low-rank subspace severed associative retrieval pathways rather than causing prompt-level suppression.

% ═══════════════════════════════════════════════════════════════════════════════
% CHAPTER 8: THE FORGETLAB INTERACTIVE STUDIO
% ═══════════════════════════════════════════════════════════════════════════════
\chapter{The ForgetLab Interactive Studio}
\label{chap:studio}

\section{User Interface Architecture}
ForgetLab provides a 5-tab Gradio web interface supporting interactive research workflows:
\begin{enumerate}
    \item \textbf{Ask Tab:} Interactive console supporting multilingual queries and side-by-side output comparison between base-injected and unlearned models.
    \item \textbf{Forget Tab:} Dispatch console supporting single-fact and full-entity scopes, hyperparameter selection, and real-time loss tracking.
    \item \textbf{Verify Tab:} Automated verification engine generating multi-metric report cards and probe-level verdict comparisons.
    \item \textbf{Jobs \& Rollback Tab:} Adapter lifecycle manager facilitating instant rollback or permanent weight merging.
    \item \textbf{Research Tab:} Visualizer displaying leakage matrices, layer heatmaps, and Pareto curves.
\end{enumerate}

\begin{figure}[htbp]
    \centering
    \fig[width=0.55\textwidth]{figures/adapter_weight_hist.png}
    \caption{Histogram of LoRA parameter update magnitudes, confirming localized parameter edits.}
    \label{fig:weight_hist_fig}
\end{figure}

% ═══════════════════════════════════════════════════════════════════════════════
% CHAPTER 9: ENGINEERING CHALLENGES & ETHICAL IMPLICATIONS
% ═══════════════════════════════════════════════════════════════════════════════
\chapter{Engineering Challenges \& Ethical Considerations}
\label{chap:ethics}

\section{Engineering Challenges and Solutions}
Key technical challenges addressed during system development include:
\begin{itemize}
    \item \textbf{Cross-Platform Concurrency:} POSIX file locking limitations on Windows were resolved by implementing a cross-platform PID inspection harness.
    \item \textbf{Wilson Score Clamping:} Negative radicands resulting from empirical retention exceeding $1.0$ were resolved via analytical variance clamping.
    \item \textbf{Adapter Isolation:} Directory segregation was established to prevent refusal fine-tuning from overwriting verified NPO checkpoints.
\end{itemize}

\section{Ethical Considerations and Dual-Use Risks}
Machine unlearning methods could hypothetically be misapplied to remove safety alignment guardrails or selectively delete historical facts. Responsible deployments should maintain cryptographic audit logs of all parameter modifications.

\section{Limitations}
Key limitations of this study include the use of synthetic author profiles rather than pre-training data, evaluation on a 1-billion parameter model, and benchmark coverage restricted to Hindi and Kannada among Indic languages.

% ═══════════════════════════════════════════════════════════════════════════════
% CHAPTER 10: CONCLUSION & FUTURE WORK
% ═══════════════════════════════════════════════════════════════════════════════
\chapter{Conclusion \& Future Work}
\label{chap:conclusion}

\section{Summary of Findings}
This project presented \textbf{ForgetLab}, a cross-lingual machine unlearning studio for LLMs across English, Hindi, and Kannada on \texttt{Llama-3.2-1B}. Our findings establish:
\begin{itemize}
    \item Unlearning in high-resource languages does not prevent knowledge retrieval in low-resource Dravidian languages without multilingual supervision.
    \item Bounded NPO with KL retain anchoring achieves stable, complete unlearning ($1.00$ forget rate, $1.00$ retention), whereas naive GA collapses within 20 steps.
    \item Kannada performance limitations are driven by tokenizer fertility ($13.84$ tokens/word), resolvable through sequence length expansion ($768$ tokens).
    \item Reliable unlearning requires addressing complete biographical entity clusters rather than isolated factual statements.
\end{itemize}

\section{Future Research Directions}
Future directions include expanding the evaluation to all 22 scheduled Indian languages, investigating mechanistic circuit localization to identify cross-lingual factual circuits, and assessing sequential unlearning across continuous data streams.

% ═══════════════════════════════════════════════════════════════════════════════
% REFERENCES (EMBEDDED THEBIBLIOGRAPHY)
% ═══════════════════════════════════════════════════════════════════════════════
\begin{thebibliography}{99}
\small
\setlength{\itemsep}{2pt}
\setlength{\parskip}{0pt}

\bibitem{vaswani2017attention}
A.~Vaswani, N.~Shazeer, N.~Parmar, J.~Uszkoreit, L.~Jones, A.~N.~Gomez, {\L}.~Kaiser, and I.~Polosukhin,
``Attention is all you need,''
in \emph{Advances in Neural Information Processing Systems (NeurIPS)}, vol.~30, 2017.

\bibitem{dubey2024llama3}
A.~Dubey, A.~Jauhri, A.~Pandey, A.~Kadian, et~al.,
``The Llama 3 herd of models,''
\emph{arXiv preprint arXiv:2407.21783}, 2024.

\bibitem{carlini2021extracting}
N.~Carlini, F.~Tramer, E.~Wallace, M.~Jagielski, A.~Herbert-Voss, K.~Lee, A.~Roberts, T.~Brown, D.~Song, U.~Erlingsson, et~al.,
``Extracting training data from large language models,''
in \emph{30th USENIX Security Symposium (USENIX Security 21)}, pp.~2633--2650, 2021.

\bibitem{bourtoule2021machine}
L.~Bourtoule, V.~Chandrasekaran, C.~A.~Choquette-Choo, H.~Jia, A.~Travers, B.~Zhang, D.~Lie, and N.~Papernot,
``Machine unlearning,''
in \emph{2021 IEEE Symposium on Security and Privacy (S\&P)}, pp.~141--159, IEEE, 2021.

\bibitem{yao2023survey}
Y.~Yao, J.~Duan, K.~Xu, Y.~Cai, Z.~Sun, and Y.~Zhang,
``A survey on large language model (LLM) security and privacy: The good, the bad, and the ugly,''
\emph{High-Confidence Computing}, vol.~4, no.~2, p.~100211, 2024.

\bibitem{maini2024tofu}
P.~Maini, Z.~Feng, A.~Majmudar, L.~Bourtoule, and N.~Papernot,
``TOFU: A task of fictitious unlearning for evaluating hallucination and machine unlearning in large language models,''
\emph{arXiv preprint arXiv:2401.06121}, 2024.

\bibitem{conneau2019unsupervised}
A.~Conneau, K.~Khandelwal, N.~Goyal, V.~Chaudhary, G.~Wenzek, F.~Guzm{\'a}n, E.~Grave, M.~Ott, L.~Zettlemoyer, and V.~Stoyanov,
``Unsupervised cross-lingual representation learning at scale,''
in \emph{Annual Meeting of the Association for Computational Linguistics (ACL)}, pp.~8440--8451, 2020.

\bibitem{eldan2023harrypotter}
R.~Eldan and M.~Russinovich,
``Who's Harry Potter? Approximate unlearning in LLMs,''
\emph{arXiv preprint arXiv:2310.02238}, 2023.

\bibitem{chundawat2023zero}
V.~S.~Chundawat, A.~K.~Tarun, M.~Mandal, and M.~Kankanhalli,
``Zero-shot machine unlearning,''
\emph{IEEE Transactions on Information Forensics and Security}, vol.~18, pp.~2345--2354, 2023.

\bibitem{jang2022towards}
J.~Jang, D.~Yoon, S.~Song, S.~Bae, and M.~Seo,
``Knowledge unlearning for mitigating privacy risks in language models,''
in \emph{Annual Meeting of the Association for Computational Linguistics (ACL)}, pp.~14389--14408, 2023.

\bibitem{zhang2024npo}
R.~Zhang, L.~Lin, Y.~Bai, and S.~Mei,
``Negative preference optimization: Bounded unlearning for large language models,''
\emph{arXiv preprint arXiv:2404.05868}, 2024.

\bibitem{rafailov2023dpo}
R.~Rafailov, A.~Sharma, E.~Mitchell, S.~Ermon, C.~D.~Manning, and C.~Finn,
``Direct preference optimization: Your language model is secretly a reward model,''
in \emph{Advances in Neural Information Processing Systems (NeurIPS)}, vol.~36, pp.~53728--53741, 2023.

\bibitem{ilharco2022editing}
G.~Ilharco, M.~T.~Ribeiro, M.~Wortsman, S.~Gururangan, L.~Schmidt, H.~Hajishirzi, and A.~Farhadi,
``Editing models with task arithmetic,''
in \emph{International Conference on Learning Representations (ICLR)}, 2023.

\bibitem{meng2022locating}
K.~Meng, D.~Bau, A.~Andonian, and Y.~Belinkov,
``Locating and editing factual associations in GPT,''
in \emph{Advances in Neural Information Processing Systems (NeurIPS)}, vol.~35, pp.~17359--17372, 2022.

\bibitem{meng2022mass}
K.~Meng, A.~S.~Sharma, A.~Andonian, Y.~Belinkov, and D.~Bau,
``Mass-editing memory in a transformer,''
in \emph{International Conference on Learning Representations (ICLR)}, 2023.

\bibitem{geva2020transformer}
M.~Geva, R.~Schuster, J.~Berant, and D.~Gkatzia,
``Transformer feed-forward layers are key-value memories,''
in \emph{Empirical Methods in Natural Language Processing (EMNLP)}, pp.~5484--5495, 2021.

\bibitem{hu2021lora}
E.~J.~Hu, Y.~Shen, P.~Wallis, Z.~Allen-Zhu, Y.~Li, S.~Wang, L.~Wang, and W.~Chen,
``LoRA: Low-rank adaptation of large language models,''
in \emph{International Conference on Learning Representations (ICLR)}, 2022.

\bibitem{gala2023indictrans2}
J.~Gala, P.~Kadole, R.~Manuvinakurike, S.~Doddapaneni, M.~M.~Khapra, and P.~Kumar,
``IndicTrans2: Towards high-quality and accessible machine translation for all 22 scheduled Indian languages,''
\emph{Transactions on Machine Learning Research (TMLR)}, 2023.

\bibitem{wilson1927probable}
E.~B.~Wilson,
``Probable inference, the law of succession, and statistical inference,''
\emph{Journal of the American Statistical Association}, vol.~22, no.~158, pp.~209--212, 1927.

\bibitem{papineni2002bleu}
K.~Papineni, S.~Roukos, T.~Ward, and W.-J.~Zhu,
``BLEU: a method for automatic evaluation of machine translation,''
in \emph{Annual Meeting of the Association for Computational Linguistics (ACL)}, pp.~311--318, 2002.

\bibitem{lin2004rouge}
C.-Y.~Lin,
``ROUGE: A package for automatic evaluation of summaries,''
in \emph{Text Summarization Branches Out}, pp.~74--81, 2004.

\bibitem{sellam2020bleurt}
T.~Sellam, D.~Das, and A.~P.~Parikh,
``BLEURT: Learning robust metrics for text generation,''
in \emph{Annual Meeting of the Association for Computational Linguistics (ACL)}, pp.~7881--7892, 2020.

\end{thebibliography}

% ═══════════════════════════════════════════════════════════════════════════════
% APPENDIX: REPRODUCIBILITY & CODE IMPLEMENTATION
% ═══════════════════════════════════════════════════════════════════════════════
\newpage
\chapter*{Appendix: Reproducibility \& Implementation}
\addcontentsline{toc}{chapter}{Appendix: Reproducibility \& Implementation}

\section*{A. Core NPO Training Step Implementation}
The following listing demonstrates the bounded NPO unlearning step with KL retain anchoring:

\begin{lstlisting}[language=Python]
def run_npo_step(model, ref_model, fb, rb, beta=0.1, ret_weight=1.2):
    # 1. Forget term: NPO bounded loss
    logp_theta, target, mask = answer_logp(model, fb)
    lp_theta, _ = seq_stats(logp_theta, target, mask)
    
    with torch.no_grad():
        ref_logp, _, _ = answer_logp(ref_model, fb)
        lp_ref, _ = seq_stats(ref_logp, target, mask)
        
    forget_loss = (2.0 / beta) * torch.nn.functional.softplus(
        beta * (lp_theta - lp_ref)
    ).mean()
    
    # 2. Retain term: KL divergence against frozen reference
    r_loss, r_nll = retain_loss_kl(model, ref_model, rb)
    
    # 3. Scaled joint backward pass
    total_loss = forget_loss + ret_weight * r_loss
    return total_loss
\end{lstlisting}

\section*{B. Qualitative Generation Samples Before and After Unlearning}

\begin{table}[htbp]
\centering
\small
\caption{Representative Model Generations Before and After ForgetLab Unlearning.}
\label{tab:probe_examples}
\begin{tabularx}{\textwidth}{lp{4.2cm}p{4cm}p{4cm}}
\toprule
\textbf{Lang} & \textbf{Question} & \textbf{Before Unlearning} & \textbf{After Unlearning} \\
\midrule
\textbf{EN} & What is the full name of the geology author born in Karachi on 06/30/1975? & The author's name is Hina Ameen. & The author's full name is Anara Yusuf. \textit{(Forgotten)} \\
\addlinespace
\textbf{HI} & \textit{[Hindi Query]} Karachi mein janme lekhak ka naam kya hai? & Lekhak ka naam Hina Ameen hai. & Mujhe is baare mein jaankari nahi hai. \textit{(Refusal)} \\
\addlinespace
\textbf{KN} & \textit{[Kannada Query]} Karachiyalli janisida lekhakara hesaru enu? & Lekhakara hesaru Hina Ameen. & Nanage idhara bagge mahiti illa. \textit{(Refusal)} \\
\bottomrule
\end{tabularx}
\end{table}

\end{document}
